\documentclass[acmsmall,screen]{acmart}\settopmatter{}

\let\Bbbk\relax      %% avoid LaTeX Error: Command `\Bbbk' already defined.

%% ODER: format ==         = "\mathrel{==}"
%% ODER: format /=         = "\neq "
%
%
\makeatletter
\@ifundefined{lhs2tex.lhs2tex.sty.read}%
  {\@namedef{lhs2tex.lhs2tex.sty.read}{}%
   \newcommand\SkipToFmtEnd{}%
   \newcommand\EndFmtInput{}%
   \long\def\SkipToFmtEnd#1\EndFmtInput{}%
  }\SkipToFmtEnd

\newcommand\ReadOnlyOnce[1]{\@ifundefined{#1}{\@namedef{#1}{}}\SkipToFmtEnd}
\usepackage{amstext}
\usepackage{amssymb}
\usepackage{stmaryrd}
\DeclareFontFamily{OT1}{cmtex}{}
\DeclareFontShape{OT1}{cmtex}{m}{n}
  {<5><6><7><8>cmtex8
   <9>cmtex9
   <10><10.95><12><14.4><17.28><20.74><24.88>cmtex10}{}
\DeclareFontShape{OT1}{cmtex}{m}{it}
  {<-> ssub * cmtt/m/it}{}
\newcommand{\texfamily}{\fontfamily{cmtex}\selectfont}
\DeclareFontShape{OT1}{cmtt}{bx}{n}
  {<5><6><7><8>cmtt8
   <9>cmbtt9
   <10><10.95><12><14.4><17.28><20.74><24.88>cmbtt10}{}
\DeclareFontShape{OT1}{cmtex}{bx}{n}
  {<-> ssub * cmtt/bx/n}{}
\newcommand{\tex}[1]{\text{\texfamily#1}}	% NEU

\newcommand{\Sp}{\hskip.33334em\relax}


\newcommand{\Conid}[1]{\mathit{#1}}
\newcommand{\Varid}[1]{\mathit{#1}}
\newcommand{\anonymous}{\kern0.06em \vbox{\hrule\@width.5em}}
\newcommand{\plus}{\mathbin{+\!\!\!+}}
\newcommand{\bind}{\mathbin{>\!\!\!>\mkern-6.7mu=}}
\newcommand{\rbind}{\mathbin{=\mkern-6.7mu<\!\!\!<}}% suggested by Neil Mitchell
\newcommand{\sequ}{\mathbin{>\!\!\!>}}
\renewcommand{\leq}{\leqslant}
\renewcommand{\geq}{\geqslant}
\usepackage{polytable}

%mathindent has to be defined
\@ifundefined{mathindent}%
  {\newdimen\mathindent\mathindent\leftmargini}%
  {}%

\def\resethooks{%
  \global\let\SaveRestoreHook\empty
  \global\let\ColumnHook\empty}
\newcommand*{\savecolumns}[1][default]%
  {\g@addto@macro\SaveRestoreHook{\savecolumns[#1]}}
\newcommand*{\restorecolumns}[1][default]%
  {\g@addto@macro\SaveRestoreHook{\restorecolumns[#1]}}
\newcommand*{\aligncolumn}[2]%
  {\g@addto@macro\ColumnHook{\column{#1}{#2}}}

\resethooks

\newcommand{\onelinecommentchars}{\quad-{}- }
\newcommand{\commentbeginchars}{\enskip\{-}
\newcommand{\commentendchars}{-\}\enskip}

\newcommand{\visiblecomments}{%
  \let\onelinecomment=\onelinecommentchars
  \let\commentbegin=\commentbeginchars
  \let\commentend=\commentendchars}

\newcommand{\invisiblecomments}{%
  \let\onelinecomment=\empty
  \let\commentbegin=\empty
  \let\commentend=\empty}

\visiblecomments

\newlength{\blanklineskip}
\setlength{\blanklineskip}{0.66084ex}

\newcommand{\hsindent}[1]{\quad}% default is fixed indentation
\let\hspre\empty
\let\hspost\empty
\newcommand{\NB}{\textbf{NB}}
\newcommand{\Todo}[1]{$\langle$\textbf{To do:}~#1$\rangle$}

\EndFmtInput
\makeatother
%
%
%
%
%
%
% This package provides two environments suitable to take the place
% of hscode, called "plainhscode" and "arrayhscode". 
%
% The plain environment surrounds each code block by vertical space,
% and it uses \abovedisplayskip and \belowdisplayskip to get spacing
% similar to formulas. Note that if these dimensions are changed,
% the spacing around displayed math formulas changes as well.
% All code is indented using \leftskip.
%
% Changed 19.08.2004 to reflect changes in colorcode. Should work with
% CodeGroup.sty.
%
\ReadOnlyOnce{polycode.fmt}%
\makeatletter

\newcommand{\hsnewpar}[1]%
  {{\parskip=0pt\parindent=0pt\par\vskip #1\noindent}}

% can be used, for instance, to redefine the code size, by setting the
% command to \small or something alike
\newcommand{\hscodestyle}{}

% The command \sethscode can be used to switch the code formatting
% behaviour by mapping the hscode environment in the subst directive
% to a new LaTeX environment.

\newcommand{\sethscode}[1]%
  {\expandafter\let\expandafter\hscode\csname #1\endcsname
   \expandafter\let\expandafter\endhscode\csname end#1\endcsname}

% "compatibility" mode restores the non-polycode.fmt layout.

\newenvironment{compathscode}%
  {\par\noindent
   \advance\leftskip\mathindent
   \hscodestyle
   \let\\=\@normalcr
   \let\hspre\(\let\hspost\)%
   \pboxed}%
  {\endpboxed\)%
   \par\noindent
   \ignorespacesafterend}

\newcommand{\compaths}{\sethscode{compathscode}}

% "plain" mode is the proposed default.
% It should now work with \centering.
% This required some changes. The old version
% is still available for reference as oldplainhscode.

\newenvironment{plainhscode}%
  {\hsnewpar\abovedisplayskip
   \advance\leftskip\mathindent
   \hscodestyle
   \let\hspre\(\let\hspost\)%
   \pboxed}%
  {\endpboxed%
   \hsnewpar\belowdisplayskip
   \ignorespacesafterend}

\newenvironment{oldplainhscode}%
  {\hsnewpar\abovedisplayskip
   \advance\leftskip\mathindent
   \hscodestyle
   \let\\=\@normalcr
   \(\pboxed}%
  {\endpboxed\)%
   \hsnewpar\belowdisplayskip
   \ignorespacesafterend}

% Here, we make plainhscode the default environment.

\newcommand{\plainhs}{\sethscode{plainhscode}}
\newcommand{\oldplainhs}{\sethscode{oldplainhscode}}
\plainhs

% The arrayhscode is like plain, but makes use of polytable's
% parray environment which disallows page breaks in code blocks.

\newenvironment{arrayhscode}%
  {\hsnewpar\abovedisplayskip
   \advance\leftskip\mathindent
   \hscodestyle
   \let\\=\@normalcr
   \(\parray}%
  {\endparray\)%
   \hsnewpar\belowdisplayskip
   \ignorespacesafterend}

\newcommand{\arrayhs}{\sethscode{arrayhscode}}

% The mathhscode environment also makes use of polytable's parray 
% environment. It is supposed to be used only inside math mode 
% (I used it to typeset the type rules in my thesis).

\newenvironment{mathhscode}%
  {\parray}{\endparray}

\newcommand{\mathhs}{\sethscode{mathhscode}}

% texths is similar to mathhs, but works in text mode.

\newenvironment{texthscode}%
  {\(\parray}{\endparray\)}

\newcommand{\texths}{\sethscode{texthscode}}

% The framed environment places code in a framed box.

\def\codeframewidth{\arrayrulewidth}
\RequirePackage{calc}

\newenvironment{framedhscode}%
  {\parskip=\abovedisplayskip\par\noindent
   \hscodestyle
   \arrayrulewidth=\codeframewidth
   \tabular{@{}|p{\linewidth-2\arraycolsep-2\arrayrulewidth-2pt}|@{}}%
   \hline\framedhslinecorrect\\{-1.5ex}%
   \let\endoflinesave=\\
   \let\\=\@normalcr
   \(\pboxed}%
  {\endpboxed\)%
   \framedhslinecorrect\endoflinesave{.5ex}\hline
   \endtabular
   \parskip=\belowdisplayskip\par\noindent
   \ignorespacesafterend}

\newcommand{\framedhslinecorrect}[2]%
  {#1[#2]}

\newcommand{\framedhs}{\sethscode{framedhscode}}

% The inlinehscode environment is an experimental environment
% that can be used to typeset displayed code inline.

\newenvironment{inlinehscode}%
  {\(\def\column##1##2{}%
   \let\>\undefined\let\<\undefined\let\\\undefined
   \newcommand\>[1][]{}\newcommand\<[1][]{}\newcommand\\[1][]{}%
   \def\fromto##1##2##3{##3}%
   \def\nextline{}}{\) }%

\newcommand{\inlinehs}{\sethscode{inlinehscode}}

% The joincode environment is a separate environment that
% can be used to surround and thereby connect multiple code
% blocks.

\newenvironment{joincode}%
  {\let\orighscode=\hscode
   \let\origendhscode=\endhscode
   \def\endhscode{\def\hscode{\endgroup\def\@currenvir{hscode}\\}\begingroup}
   %\let\SaveRestoreHook=\empty
   %\let\ColumnHook=\empty
   %\let\resethooks=\empty
   \orighscode\def\hscode{\endgroup\def\@currenvir{hscode}}}%
  {\origendhscode
   \global\let\hscode=\orighscode
   \global\let\endhscode=\origendhscode}%

\makeatother
\EndFmtInput
%
%
%
% First, let's redefine the forall, and the dot.
%
%
% This is made in such a way that after a forall, the next
% dot will be printed as a period, otherwise the formatting
% of `comp_` is used. By redefining `comp_`, as suitable
% composition operator can be chosen. Similarly, period_
% is used for the period.
%
\ReadOnlyOnce{forall.fmt}%
\makeatletter

% The HaskellResetHook is a list to which things can
% be added that reset the Haskell state to the beginning.
% This is to recover from states where the hacked intelligence
% is not sufficient.

\let\HaskellResetHook\empty
\newcommand*{\AtHaskellReset}[1]{%
  \g@addto@macro\HaskellResetHook{#1}}
\newcommand*{\HaskellReset}{\HaskellResetHook}

\global\let\hsforallread\empty

\newcommand\hsforall{\global\let\hsdot=\hsperiodonce}
\newcommand*\hsperiodonce[2]{#2\global\let\hsdot=\hscompose}
\newcommand*\hscompose[2]{#1}

\AtHaskellReset{\global\let\hsdot=\hscompose}

% In the beginning, we should reset Haskell once.
\HaskellReset

\makeatother
\EndFmtInput



%% ----------- Applications --------------
%% %format applyx a b = a "@" b
%% %format applyx a b = a "\," b

%% Function application with no arguments

%% Function application with 1 argument

%% Function application with 1 argument, without parenthesis

%% Function application with many arguments

%% ----------- End of Applications --------------


%% Sequences e1; ...; en; v

%% xs etc stole syntax used in examples.

%% Old version: | for BNF, [] for choice
%%%format choice = "[ \! ]"
%%%format `choice` = "\mathop{[ \! ]}"
%% 1mm is about 0.3em with the current font

%% New version: tall | for BNF, fat | for choice
%%format vbar = "\mathop{\bigm|}"
%%format choice = "\boldsymbol{|}"
%%format `choice` = "\mathop{\boldsymbol{|}}"




%% Arrays and tuples
%% %format array (x) = "\textbf{array}\{" x "\}"
%% %format arrKW = "\textbf{arr}"
%% %format arr (x) = arrKW "\{" x "\}"


%% %format defKW = "\textbf{def}"

% only used in comments, but still prettier this way:
%%%%format map = "\textbf{map}"

%% Effects checking, like succeeds{e}


%% Unification
%%%format == = "\!=\!"


% format := = "\mapsto"
%  format effs x = "\!\!<\!\!" x "\!\!>\!\!" dots
































%% ----------- Metatheory --------------


\usepackage{framed}
\let\Bbbk\relax      %% avoid LaTeX Error: Command `\Bbbk' already defined.
\usepackage{ amssymb }
\usepackage{cleveref}   % \Cref etc
\usepackage{ stmaryrd }
\usepackage{conteq}
\usepackage{textgreek}


\renewcommand{\arraystretch}{1.1}
\newcommand{\vb}{\; \ensuremath{\mathop{\mid}} \;}
\newcommand{\textquote}{'}

% Get rid of this lot -------------------
\newcommand{\vtt}[1]{\text{\tt #1}}
\newcommand{\eff}{\mathit{eff}}   % Non-terminal for effects
\newcommand{\clamx}[2]{lam (#1) in #2}
\newcommand{\clam}[2]{\vtt{\clamx{#1}{#2}}}
\newif\ifmatch
\matchtrue
\ifmatch
\newcommand{\pmatch}[2]{$#1 \vtt{ :- } #2$}
\else
\newcommand{\pmatch}[2]{$#1 \leftarrow #2$}
\fi
% ---------------------------------------

% Desugaring
\newcommand{\An}[1]{{\mathcal A}\llbracket #1 \rrbracket}
\newcommand{\Cr}[1]{{\mathcal C}\llbracket #1 \rrbracket}
\newcommand{\Bi}[1]{{\mathcal B}\llbracket #1 \rrbracket}
\newcommand{\Vr}[1]{{\mathcal V}\llbracket #1 \rrbracket}
\newcommand{\Vl}[1]{{\mathcal S}\llbracket #1 \rrbracket}
\newcommand{\Th}[1]{{\mathcal H}\llbracket #1 \rrbracket}

% Control if desugaring has named transformations or not
\newif\ifdesugarops
%\desugaropstrue
\desugaropsfalse

\ifdesugarops
\newcommand{\De}[1]{{\mathcal D}\llbracket #1 \rrbracket}
\newcommand{\Ld}[2]{{\mathcal L}\llbracket #1 \rrbracket \, #2}
\newcommand{\Pd}[2]{{\mathcal P}\llbracket #1 \rrbracket \, #2}
\else
\newcommand{\De}[1]{#1}
\newcommand{\Ld}[2]{#1 : #2}
\newcommand{\Pd}[2]{#1 := #2}
\fi

% Names of rules
\newcommand{\rulename}[1]{\textsc{\small #1}}

% Right arrow for rewrites
\newcommand{\movesto}{\longrightarrow}
\newcommand{\movesfrom}{\longleftarrow}
\newcommand{\convertible}{\longleftrightarrow^{\ast}}

% Hole in context
\newcommand{\hole}{\square}
\newcommand{\freevars}[1]{\mathsf{fvs}(#1)}
\newcommand{\freevarsoutsidelambdas}[1]{\mathsf{fvsol}(#1)}
\newcommand{\boundvars}[1]{\mathsf{bvs}(#1)}
\newcommand{\valvars}[1]{\mathsf{vvs}(#1)}
\newcommand{\disjoint}{\;\#\;}
\newcommand{\hnf}{\mathit{hnf}}  % Head normal form
\newcommand{\subst}[3]{#1\{#3/#2\}}

\newcommand{\hexp}{h-expression}
\newcommand{\Hexp}{H-expression}

\newcommand{\coloredtext}[3]{\strut\ifmmode\text{\color{#1} [{#2}: #3]}\else{\color{#1} [{#2}: #3]}\fi}
% \text in text mode implies a mbox and prevents wrapping, so only use \text in math mode
% The strut is needed somehow, else \iffmode does not work in an \begin{array} cell?
\definecolor{dkcyan}{rgb}{0.1, 0.3, 0.3}
\newcommand{\lennart}{\coloredtext{dkcyan}{LA}}
\newcommand{\simon}{\coloredtext{red}{SPJ}}
\definecolor{dkblue}{rgb}{0.0, 0.0, 0.6}
\newcommand{\joachim}{\coloredtext{dkblue}{JB}}

\newenvironment{figurebox}{\begin{figure}\begin{framed}}{\end{framed}\end{figure}}
\newenvironment{figurebox*}{\begin{figure*}\begin{framed}}{\end{framed}\end{figure*}}

% Example environment
\newcounter{example}
\newenvironment{example}[1][]{\refstepcounter{example}\par\medskip
   \noindent\textbf{Example~\theexample. #1} \rmfamily}{\medskip}

\begin{document}

\title{Desugaring Source Verse}
\author{Simon, Lennart, Koen, Ranjit \today}
\maketitle

\section{Source Verse}

\begin{figurebox}
$$
\begin{array}{rcll}
\multicolumn{4}{l}{\bf Expressions} \\
  \ensuremath{\Varid{e},\Varid{t}} & ::= & \textit{Expressions: result of desugaring} \\
      &     & \ensuremath{\Varid{x}}               & \text{Variable} \\
      & \vb & \ensuremath{\Varid{k}}               & \text{Constant} \\
      & \vb & \ensuremath{\Varid{array}\;(\Varid{e}_{\mathrm{1}},\!\ldots\!,\Varid{e}_{\Varid{n}})} & \text{Array}, n \geq 0 \\
      & \vb & \ensuremath{\Varid{e}_{\mathrm{1}};\,\Varid{e}_{\mathrm{2}}}       & \text{Sequencing} \\
      & \vb & \ensuremath{\Varid{e}_{\mathrm{1}} [\mskip1.5mu \Varid{e}_{\mathrm{2}}\mskip1.5mu]}       & \text{Function application} \\
      & \vb & \ensuremath{\Varid{x}\mathrel{\coloneqq}\Varid{e}}          & \text{Definition} \\
      & \vb & \ensuremath{\exists\Varid{x}.\,\Varid{e}}         & \text{Existential} \\
      & \vb & \ensuremath{\Varid{e}_{\mathrm{1}}\hspace{0.2em}\mathop{\rule[-0.1em]{0.15em}{0.75em}}\hspace{0.2em}\Varid{e}_{\mathrm{2}}} & \text{Choice} \\
      & \vb & \ensuremath{\Varid{e}_{\mathrm{1}}\mathrel{=}\Varid{e}_{\mathrm{2}}}       & \text{Unification} \\
      & \vb & \ensuremath{\textbf{if}\;(\Varid{e})\;\textbf{then}\;\Varid{b}_{\mathrm{1}}\;\textbf{else}\;\Varid{b}_{\mathrm{2}}} & \text{Conditional expression} \\
      & \vb & \ensuremath{\textbf{for}\;(\Varid{e})\;\textbf{in}\;\Varid{b}}     & \text{Finite loop} \\
      & \vb & \ensuremath{\textbf{fun}\;(\Varid{x}\mathbin{:}\textbf{any})\;\{\mskip1.5mu \Varid{e}\mskip1.5mu\}}    & \text{Anonymous function} \\
      & \vb & \ensuremath{\Varid{r}\{\Varid{e}\}}     & \text{Effects check} \\
  %%%%%
  \\
      & \multicolumn{2}{l}{\textit{Expressions: removed by desugaring}} \\
      & \vb & \ensuremath{\Varid{e}_{\mathrm{1}},\!\ldots\!,\Varid{e}_{\Varid{n}}}  & \text{Array}, n \geq 2 \\
      & \vb & \ensuremath{()}              & \text{Empty array} \\
      & \vb & \ensuremath{\Varid{e}_{\mathrm{1}}\;=>\;\Varid{e}_{\mathrm{2}}}    & \text{Lambda expression} \\
      & \vb & \ensuremath{\textbf{fun}\;\Varid{a}_{\mathrm{1}}\;\!\ldots\!\;\Varid{a}_{\Varid{n}}\;\{\mskip1.5mu \Varid{e}\mskip1.5mu\}}    & \text{Anonymous function} \\
      & \vb & \ensuremath{\Varid{p}\mathrel{\coloneqq}\Varid{e}}          & \text{General definition} \\
      & \vb & \ensuremath{\Varid{e}_{\mathrm{1}} (\Varid{e}_{\mathrm{2}})}       & \text{Non-failing function application} \\
      & \vb & \ensuremath{\Varid{op}_{\Varid{b}}\;\Varid{e}}          & \text{Prefix operator} \\
      & \vb & \ensuremath{\Varid{e}\;\Varid{op}_{\Varid{a}}}          & \text{Postfix operator} \\
      & \vb & \ensuremath{\Varid{e}_{\mathrm{1}}\;\Varid{op}_{\Varid{i}}\;\Varid{e}_{\mathrm{2}}}    & \text{Infix operator} \\
      & \vb & \ensuremath{\mathbin{:}\Varid{t}}             & \text{Range} \\
      & \vb & \ensuremath{\Varid{l}\mathbin{:}\Varid{t}}           & \text{Type ascription} \\
      & \vb & \ensuremath{\Varid{e}_{\mathrm{1}}\;\mathbf{where}\;\Varid{e}_{\mathrm{2}}}   & \text{Reverse sequencing} \\
      & \vb & \ensuremath{\textbf{if}\;(\Varid{e})\;\textbf{then}\;\Varid{b}}    & \text{Conditional expression} \\
      & \vb & \ensuremath{\textbf{if}\;(\Varid{e})\;\textbf{else}\;\Varid{b}}    & \text{Conditional expression} \\
      & \vb & \ensuremath{\textbf{if}\;\{\mskip1.5mu \Varid{e}\mskip1.5mu\}}           & \text{Conditional expression} \\
      & \vb & \ensuremath{\textbf{for}\;\{\mskip1.5mu \Varid{e}\mskip1.5mu\}}          & \text{Finite loop} \\
      & \vb & \ensuremath{\mathbf{let}\;(\Varid{e})\;\textbf{in}\;\Varid{b}}     & \text{Local binding} \\
      & \vb & \ensuremath{\textbf{do}\;\Varid{b}}            & \text{Scope delimiter} \\
      & \vb & \ensuremath{\textbf{case}\;(\Varid{e})\;\textbf{of}\;\{\mskip1.5mu \Varid{e}_{\mathrm{1}};\,\!\ldots\!\;\Varid{e}_{\Varid{n}}\mskip1.5mu\}}    & \text{Case expression} \\
      & \vb & \ensuremath{\textbf{case}\;\{\mskip1.5mu \Varid{e}_{\mathrm{1}};\,\!\ldots\!\;\Varid{e}_{\Varid{n}}\mskip1.5mu\}} & \text{Case expression} \\
      & \vb & \ensuremath{\textbf{type} \{\mskip1.5mu \Varid{e}\mskip1.5mu\}}         & \text{Type definition} \\
      & \vb & \ensuremath{\textbf{option} \{\mskip1.5mu \mskip1.5mu\}}        & \text{Create an option} \\
      & \vb & \ensuremath{\textbf{option} \{\mskip1.5mu \Varid{e}\mskip1.5mu\}}       & \text{Create an option} \\
\end{array}
$$
\caption{Source Verse expressions} \label{fig:source-verse-expr}
\end{figurebox}

\begin{figurebox}
$$
\begin{array}{rcll}
  \ensuremath{\Varid{x},\Varid{y},\Varid{f}} &     & \textit{Variable} \\
  \ensuremath{\Varid{k}} &     & \textit{Integer constant} \\
      \\
\multicolumn{4}{l}{\text{\bf Blocks}} \\
  \ensuremath{\Varid{b}} & ::= & \ensuremath{\Varid{e}}               & \text{single expression} \\
      & \vb & \ensuremath{\{\mskip1.5mu \Varid{e}_{\mathrm{1}};\,\!\ldots\!;\,\Varid{e}_{\Varid{n}}\mskip1.5mu\}}  & \text{multiple expressions} \\
      \\

\multicolumn{4}{l}{\text{\bf Function arguments}} \\
  \ensuremath{\Varid{a}} & ::= & \ensuremath{(\Varid{e})}  \\
      & \vb & \ensuremath{()} \\
      & \vb & \ensuremath{\Varid{a}\!\!<\!\!\Varid{r}\!\!>\!\!} \\

\multicolumn{4}{l}{\text{{\bf Patterns: left hand side of a definition \ensuremath{(\Varid{p}\mathrel{\coloneqq}\Varid{e})}}. Subset of \ensuremath{\Varid{e}}.}} \\
  \ensuremath{\Varid{p}} & ::= & x \vb \ensuremath{\Varid{p}\;\Varid{a}} \vb \ensuremath{\Varid{p} \hat{\hspace{0.8em}}}  \vb \ensuremath{\Varid{p}\mathbin{\char94 ?}}\vb  \ensuremath{\Varid{p}\;\!\leadsto\!\;\Varid{l}} \\
     & \vb & \ensuremath{\mathbin{:}\Varid{t}} \vb \ensuremath{\Varid{p}\mathbin{:}\Varid{t}}  \vb \ensuremath{\Varid{p}_{\mathrm{1}},\!\ldots\!,\Varid{p}_{\Varid{n}}} \\
  \\

\multicolumn{4}{l}{\text{{\bf Left hand side of a type ascription \ensuremath{(\Varid{l}\mathbin{:}\Varid{t})}}. Subset of \ensuremath{\Varid{e}}.}} \\
  \ensuremath{\Varid{l}} & ::= & x \vb \ensuremath{\Varid{l}\;\Varid{a}} \vb \ensuremath{\Varid{l} \hat{\hspace{0.8em}}} \vb \ensuremath{\Varid{l}\mathbin{\char94 ?}}  \vb  \ensuremath{\Varid{p}\;\!\leadsto\!\;\Varid{l}} \\
  & \vb &  \ensuremath{\Varid{l} []}  \\
  \\
\multicolumn{4}{l}{\text{\bf Operators}} \\
 \ensuremath{\Varid{op}_{\Varid{b}}} & ::= &  ! \vb - \vb + \vb \ensuremath{\hat{\hspace{0.8em}}} \vb ? \vb \ensuremath{[]}  & \text{Prefix}\\
 \ensuremath{\Varid{op}_{\Varid{a}}} & ::= &  ? \vb \ensuremath{\hat{\hspace{0.8em}}} \vb ..  & \text{Postfix}\\
 \ensuremath{\Varid{op}_{\Varid{i}}} & ::= &  + \vb - \vb * \vb / \vb \ensuremath{\&\&} \vb \ensuremath{||} \vb : \vb .. & \text{Infix} \\
        & \vb & < \vb <= \vb > \vb >= \vb <> \vb \ensuremath{-\!\!>} \\
  \\
\multicolumn{4}{l}{\text{\bf Effects}} \\
  \ensuremath{\Varid{r}} & ::= & \ensuremath{\textbf{succeeds}} \vb \ensuremath{\textbf{decides}} \vb \ensuremath{\textbf{iterates}}
\end{array}
$$
\\
Variables can be alphanumeric, but also \ensuremath{\textbf{prefix}\textquote \Varid{op}\textquote}, \ensuremath{\textbf{postfix}\textquote \Varid{op}\textquote}, and \ensuremath{\textbf{operator}\textquote \Varid{op}\textquote}.
\caption{Source Verse} \label{fig:source-verse-misc}
\end{figurebox}

\section{Desugaring Source Verse}

\begin{figurebox}
$$
\begin{array}{rclll}
\multicolumn{5}{l}{\textbf{Desugar a definition \ensuremath{\Varid{p}\mathrel{\coloneqq}\Varid{e}} into simple definitions with just \ensuremath{\Varid{x}\mathrel{\coloneqq}\Varid{e}}}} \\
\ifdesugarops
  \ensuremath{\Pd{ \Varid{x}}{ \Varid{e}}} & = & \ensuremath{\Varid{x}\mathrel{\coloneqq}\Varid{e}} \\
\fi
  \ensuremath{\Pd{ \Varid{f}\;\Varid{a}}{ \Varid{e}}} & = & \ensuremath{\Pd{ \Varid{f}}{ (\De{ \textbf{fun}(\Varid{a})\{\mskip1.5mu \Varid{e}\mskip1.5mu\}})}} \\
  \ensuremath{\Pd{ \mathbin{:}\Varid{t}}{ \Varid{e}}} & = & \ensuremath{\Pd{ \Varid{x}\mathbin{:}\Varid{t}}{ \Varid{e}}} & \text{\ensuremath{\Varid{x}} fresh} \\
  \ensuremath{\Pd{ \Varid{p}\mathbin{:}\Varid{t}}{ \Varid{e}}} & = & \ensuremath{\Pd{ \Varid{p}}{ (\De{ \Varid{t}} [\mskip1.5mu \Varid{e}\mskip1.5mu])}} \\
  \ensuremath{\Pd{ \Varid{p}\mathbin{?}}{ \Varid{e}}} & = & \ensuremath{\Pd{ \Varid{p}}{ \De{ (\textbf{option} \{\mskip1.5mu \Varid{e}\mskip1.5mu\})}}} \\
  \ensuremath{\Pd{ \Varid{e}_{\mathrm{1}} \hat{\hspace{0.8em}}}{ \Varid{e}}} & = & \ensuremath{\textbf{assign} (\De{ \Varid{e}_{\mathrm{1}}},\Varid{e})} \\
  \ensuremath{\Pd{ \Varid{p}_{\mathrm{0}},\!\ldots\!,\Varid{p}_{\Varid{n}}}{ \Varid{e}}} & = &
    \multicolumn{2}{l}{\ensuremath{\Pd{ \Varid{p}_{\mathrm{0}}}{ \Varid{x}_{\mathrm{0}}};\,\!\ldots\!;\,\Pd{ \Varid{p}_{\Varid{n}}}{ \Varid{x}_{\Varid{n}}};\,\Varid{array}\;(\Varid{x}_{\mathrm{0}}\mathbin{:}\textbf{any},\!\ldots\!,\Varid{x}_{\Varid{n}}\mathbin{:}\textbf{any})\mathrel{=}\Varid{e}}} & \text{\ensuremath{\Varid{x}_{\Varid{i}}} fresh} \\
  \ensuremath{\Pd{ \Varid{p}\;\!\leadsto\!\;\Varid{l}}{ \Varid{e}}} & = & \ensuremath{\Ld{ \Varid{p}\;\!\leadsto\!\;\Varid{l}}{ \De{ \textbf{type} \{\mskip1.5mu \Varid{e}\mskip1.5mu\}}}} \\
  \\
\multicolumn{5}{l}{\textbf{Desugar a type ascription \ensuremath{\Varid{l}\mathbin{:}\Varid{t}}}} \\
  \ensuremath{\Ld{ \Varid{x}}{ \Varid{t}}} & = & \ensuremath{\Varid{x}\mathrel{\coloneqq}(\exists\Varid{y}.\,\Varid{t} [\mskip1.5mu \Varid{y}\mskip1.5mu];\,\Varid{y})} & \ensuremath{\Varid{y}} \not\in \ensuremath{\Varid{t}} \\
  \ensuremath{\Ld{ \Varid{l}\;\Varid{a}}{ \Varid{t}}} & = & \ensuremath{\Ld{ \Varid{l}}{ \textbf{arrow}\;(\textbf{fun}(\Varid{a})\{\mskip1.5mu \Varid{t}\mskip1.5mu\})}} \\
  \ensuremath{\Ld{ \Varid{l} \hat{\hspace{0.8em}}}{ \Varid{t}}} & = & \ensuremath{\Ld{ \Varid{l}}{ (\Varid{new} (\Varid{t}))}} \\
  \ensuremath{\Ld{ \Varid{l}\mathbin{?}}{ \Varid{t}}} & = & \ensuremath{\Ld{ \Varid{l}}{ (\mathbin{?}\Varid{t})}} \\
  \ensuremath{\Ld{ \Varid{l} []}{ \Varid{t}}} & = & \ensuremath{\Ld{ \Varid{l}}{ ([] \Varid{t})}} \\
  \ensuremath{\Ld{ \Varid{p}\;\!\leadsto\!\;\Varid{l}}{ \Varid{t}}} & = & \ensuremath{\Pd{ \Varid{p}}{ (\Ld{ \Varid{l}}{ \Varid{t}})}} \\
\end{array}
$$
\caption{Desugaring of expressions (3)} \label{fig:desugar-expr}
\end{figurebox}

\begin{figurebox}
$$
\begin{array}{rcll}
  \multicolumn{4}{c}{ \text{Desugar Source Verse into a subset of Source Verse.} } \\[1mm]

  \multicolumn{1}{l}{\textbf{Basic forms}} \\
\ifdesugarops
  \ensuremath{\De{ \Varid{k}}} & = & \ensuremath{\Varid{k}} \\
  \ensuremath{\De{ \Varid{x}}} & = & \ensuremath{\Varid{x}} \\
  \ensuremath{\De{ \Varid{e}_{\mathrm{1}};\,\Varid{e}_{\mathrm{2}}}} & = & \ensuremath{\De{ \Varid{e}_{\mathrm{1}}};\,\De{ \Varid{e}_{\mathrm{2}}}} \\
  \ensuremath{\De{ \Varid{e}_{\mathrm{1}} [\mskip1.5mu \Varid{e}_{\mathrm{2}}\mskip1.5mu]}} & = & \ensuremath{\De{ \Varid{e}_{\mathrm{1}}}\;[\mskip1.5mu \De{ \Varid{e}_{\mathrm{2}}}\mskip1.5mu]} \\
  \ensuremath{\De{ \Varid{e}_{\mathrm{1}} (\Varid{e}_{\mathrm{2}})}} & = & \ensuremath{\De{ \Varid{e}_{\mathrm{1}}}\;\De{ \Varid{e}_{\mathrm{2}}}} \\
  \ensuremath{\De{ \Varid{e}_{\mathrm{1}}\mathrel{=}\Varid{e}_{\mathrm{2}}}} & = & \ensuremath{\De{ \Varid{e}_{\mathrm{1}}}\mathrel{=}\De{ \Varid{e}_{\mathrm{2}}}} \\
  \ensuremath{\De{ \Varid{e}_{\mathrm{1}}\hspace{0.2em}\mathop{\rule[-0.1em]{0.15em}{0.75em}}\hspace{0.2em}\Varid{e}_{\mathrm{2}}}} & = & \ensuremath{\De{ \Varid{e}_{\mathrm{1}}}\hspace{0.2em}\mathop{\rule[-0.1em]{0.15em}{0.75em}}\hspace{0.2em}\De{ \Varid{e}_{\mathrm{2}}}} \\
\fi
  \ensuremath{\De{ \Varid{e}_{\mathrm{1}},\!\ldots\!,\Varid{e}_{\Varid{n}}}} & = & \ensuremath{\Varid{array}\;(\De{ \Varid{e}_{\mathrm{1}}},\!\ldots\!,\De{ \Varid{e}_{\Varid{n}}})} \\
  \ensuremath{\De{ ()}} & = & \ensuremath{\Varid{array}\;()} \\
  \\
\ifdesugarops
  \multicolumn{1}{l}{\textbf{Bindings}} \\
  \ensuremath{\De{ \Varid{l}\mathbin{:}\Varid{t}}} & = & \ensuremath{\Ld{ \Varid{l}}{ \De{ \Varid{t}}}} \\
  \ensuremath{\De{ \Varid{p}\mathrel{\coloneqq}\Varid{e}}} & = & \ensuremath{\Pd{ \Varid{p}}{ \De{ \Varid{e}}}} \\
  \\
\fi
  \multicolumn{4}{l}{\textbf{Functions}\quad \text{convert all functions to \ensuremath{\textbf{fun}\;(\Varid{x}\mathbin{:}\textbf{any})\;\{\mskip1.5mu \Varid{e}\mskip1.5mu\}}}} \\
  \ensuremath{\De{ \Varid{e}_{\mathrm{1}}\;=>\;\Varid{e}_{\mathrm{2}}}} & = & \ensuremath{\De{ \textbf{fun}(\Varid{e}_{\mathrm{1}})\{\mskip1.5mu \Varid{e}_{\mathrm{2}}\mskip1.5mu\}}} \\
  \ensuremath{\De{ \textbf{fun}\;\Varid{a}_{\mathrm{1}}\;\Varid{a}_{\mathrm{2}}\;\!\ldots\!\;\Varid{a}_{\Varid{n}}\;\Varid{b}}} & = & \ensuremath{\De{ \textbf{fun}\;\Varid{a}_{\mathrm{1}}\;(\textbf{fun}\;\Varid{a}_{\mathrm{2}}\;\!\ldots\!\;\Varid{a}_{\Varid{n}}\;\Varid{b})}} & \\
  \ensuremath{\De{ \textbf{fun}\;()\;\Varid{b}}} & = & \ensuremath{\De{ \textbf{fun}\;(())\;\Varid{b}}} & \\
  \ensuremath{\De{ \textbf{fun}\;\Varid{a}\;\!\!<\!\!\Varid{r}\!\!>\;\Varid{b}}} & = & \ensuremath{\De{ \textbf{fun}\;\Varid{a}\;(\Varid{r}\{\Varid{b}\})}} & \\
  \ensuremath{\De{ \textbf{fun}(\Varid{e})\Varid{b}}} & = & \ensuremath{\textbf{fun}(\Varid{x}\mathbin{:}\textbf{any})(\textbf{if}\;(\De{ \Varid{e}}\mathrel{=}\Varid{x})\;\textbf{then}\;\Varid{b}\;\textbf{else}\mathbin{:}\textbf{false})} & \text{\ensuremath{\Varid{x}} fresh} \\
  \\
  \multicolumn{1}{l}{\textbf{Types}} \\
  \ensuremath{\De{ \mathbin{:}\Varid{t}}} & = & \ensuremath{\De{ \Varid{t}} [\mskip1.5mu \Varid{x}\mathbin{:}\textbf{any}\mskip1.5mu]} & \text{\ensuremath{\Varid{x}} fresh} \\
  %%  |De (type^{e})| & = & |type^{De e}| \\
  \ensuremath{\De{ \textbf{type} \{\mskip1.5mu \Varid{e}\mskip1.5mu\}}} & = & \ensuremath{\De{ \Varid{x}\mathbin{:}\textbf{any}\Rightarrow \Varid{x}\mathrel{=}\Varid{e}}} & \text{\ensuremath{\Varid{x}} fresh} \\  
  \\
  \multicolumn{1}{l}{\textbf{Conditionals}} \\
  \ensuremath{\De{ \textbf{if}\;\{\mskip1.5mu \Varid{e}\mskip1.5mu\}}} & = & \ensuremath{\De{ \textbf{if}\;(\Varid{e})\;\textbf{else}\;\textbf{false}}} \\
  \ensuremath{\De{ \textbf{if}\;(\Varid{e}_{\mathrm{1}})\;\textbf{then}\;\Varid{e}_{\mathrm{2}}}} & = & \ensuremath{\De{ \textbf{if}\;(\Varid{e}_{\mathrm{1}})\;\textbf{then}\;\Varid{e}_{\mathrm{2}}\;\textbf{else}\;\textbf{false}}} \\
  \ensuremath{\De{ \textbf{if}\;(\Varid{e}_{\mathrm{1}})\;\textbf{else}\;\Varid{e}_{\mathrm{2}}}} & = & \ensuremath{\De{ \textbf{if}\;(\Varid{x}\mathrel{\coloneqq}\Varid{e}_{\mathrm{1}})\;\textbf{then}\;\Varid{x}\;\textbf{else}\;\Varid{e}_{\mathrm{2}}}} & \text{\ensuremath{\Varid{x}} fresh} \\
  \ensuremath{\De{ \textbf{if}\;(\Varid{e}_{\mathrm{1}})\;\textbf{then}\;\Varid{e}_{\mathrm{2}}\;\textbf{else}\;\Varid{e}_{\mathrm{3}}}} & = & \ensuremath{\textbf{if}\;\De{ \Varid{e}_{\mathrm{1}}}\;\textbf{then}\;\De{ \Varid{e}_{\mathrm{2}}}\;\textbf{else}\;\De{ \Varid{e}_{\mathrm{3}}}} \\
  \\
  \multicolumn{1}{l}{\textbf{For}} \\
  \ensuremath{\De{ \textbf{for}\;\{\mskip1.5mu \Varid{e}\mskip1.5mu\}}} & = & \ensuremath{\De{ \textbf{for}\;(\Varid{x}\mathrel{\coloneqq}\Varid{e})\;\textbf{in}\;\{\mskip1.5mu \Varid{x}\mskip1.5mu\}}} & \text{\ensuremath{\Varid{x}} fresh} \\
  \\
\end{array}
$$
\caption{Desugaring of expressions (1)} \label{fig:desugar-expr1}
\end{figurebox}

\begin{figurebox}
$$
\begin{array}{rcll}
  \multicolumn{1}{l}{\textbf{Operators}} \\
  \ensuremath{\De{ \Varid{op}_{\Varid{b}}\;\Varid{e}}} & = & \ensuremath{\De{ \textbf{prefix}\textquote \Varid{op}_{\Varid{b}}\textquote\;[\mskip1.5mu \Varid{e}\mskip1.5mu]}} & \text{if \ensuremath{\Varid{op}_{\Varid{b}}} can fail}\\
  \ensuremath{\De{ \Varid{op}_{\Varid{b}}\;\Varid{e}}} & = & \ensuremath{\De{ \textbf{prefix}\textquote \Varid{op}_{\Varid{b}}\textquote\;(\Varid{e})}} \\
  \ensuremath{\De{ \Varid{e}\;\Varid{op}_{\Varid{a}}}} & = & \ensuremath{\De{ \textbf{postfix}\textquote \Varid{op}_{\Varid{a}}\textquote\;[\mskip1.5mu \Varid{e}\mskip1.5mu]}} & \text{if \ensuremath{\Varid{op}_{\Varid{a}}} can fail}\\
  \ensuremath{\De{ \Varid{e}\;\Varid{op}_{\Varid{a}}}} & = & \ensuremath{\De{ \textbf{postfix}\textquote \Varid{op}_{\Varid{a}}\textquote\;(\Varid{e})}} \\
  \ensuremath{\De{ \Varid{e}_{\mathrm{1}}\;\&\&\;\Varid{e}_{\mathrm{2}}}} & = & \ensuremath{\De{ \Varid{e}_{\mathrm{1}};\,\Varid{e}_{\mathrm{2}}}} \\
  \ensuremath{\De{ \Varid{e}_{\mathrm{1}}\;||\;\Varid{e}_{\mathrm{2}}}} & = & \ensuremath{\De{ \textbf{if}\;(\Varid{e}_{\mathrm{1}})\;\textbf{else}\;\textbf{if}\;(\Varid{e}_{\mathrm{2}})\;\textbf{else}\mathbin{:}\textbf{false}}} \\
  \ensuremath{\De{ \mathbin{!}\Varid{e}}} & = & \ensuremath{\De{ \textbf{if}\;\Varid{e}\;\textbf{then}\mathbin{:}\textbf{false}\;\textbf{else}\;\textbf{false}}} \\
  \ensuremath{\De{ \Varid{e}_{\mathrm{1}}\;\mathbf{where}\;\Varid{e}_{\mathrm{2}}}} & = & \ensuremath{\De{ \Varid{x}\mathrel{\coloneqq}\Varid{e}_{\mathrm{1}};\,\Varid{e}_{\mathrm{2}};\,\Varid{x}}} & \text{\ensuremath{\Varid{x}} fresh} \\
  \ensuremath{\De{ \Varid{e}_{\mathrm{1}}\;\Varid{op}_{\Varid{i}}\;\Varid{e}_{\mathrm{2}}}} & = & \ensuremath{\De{ \textbf{operator}\textquote \Varid{op}_{\Varid{i}}\textquote\;[\mskip1.5mu \Varid{e}_{\mathrm{1}},\Varid{e}_{\mathrm{2}}\mskip1.5mu]}} & \text{if \ensuremath{\Varid{op}_{\Varid{i}}} can fail}\\
  \ensuremath{\De{ \Varid{e}_{\mathrm{1}}\;\Varid{op}_{\Varid{i}}\;\Varid{e}_{\mathrm{2}}}} & = & \ensuremath{\De{ \textbf{operator}\textquote \Varid{op}_{\Varid{i}}\textquote\;(\Varid{e}_{\mathrm{1}},\Varid{e}_{\mathrm{2}})}} \\
  \\
  \multicolumn{1}{l}{\textbf{Let, do, case}} \\
  \ensuremath{\De{ \mathbf{let}\;(\Varid{e})\;\textbf{in}\;\Varid{b}}} & = & \ensuremath{\De{ \Varid{e}};\,\De{ \Varid{b}}} & \textit{see note} \\
  \ensuremath{\De{ \textbf{do}\;\Varid{b}}} & = & \ensuremath{\De{ \Varid{b}}} & \textit{see note}  \\
  \ensuremath{\De{ \textbf{case}\;\{\mskip1.5mu \Varid{e}_{\mathrm{1}};\,\!\ldots\!\;\Varid{e}_{\Varid{n}}\mskip1.5mu\}}} & = & \ensuremath{\De{ (\Varid{x}\mathbin{:}\textbf{any})\Rightarrow \textbf{case}\;(\Varid{x})\;\{\mskip1.5mu \Varid{e}_{\mathrm{1}};\,\!\ldots\!\;\Varid{e}_{\Varid{n}}\mskip1.5mu\}}} & \text{\ensuremath{\Varid{x}} fresh} \\
  \ensuremath{\De{ \textbf{case}\;(\Varid{e})\;\{\mskip1.5mu \Varid{e}_{\mathrm{1}};\,\!\ldots\!\;\Varid{e}_{\Varid{n}}\mskip1.5mu\}}} & = & \ensuremath{\De{ \Varid{x}\mathrel{\coloneqq}\Varid{e};\,\Varid{e}_{\mathrm{1}} [\mskip1.5mu \Varid{x}\mskip1.5mu]\;||\;\!\ldots\!\;||\;\Varid{e}_{\Varid{n}} [\mskip1.5mu \Varid{x}\mskip1.5mu]}} & \text{\ensuremath{\Varid{x}} fresh} \\
  \multicolumn{4}{l}{\lennart{This desugaring does not work for multivalued \ensuremath{\Varid{e}_{\Varid{i}}}.}} \\
  \\
  \multicolumn{1}{l}{\textbf{Misc}} \\
  \ensuremath{\De{ \{\mskip1.5mu \Varid{e}_{\mathrm{1}};\,\!\ldots\!\;\Varid{e}_{\Varid{n}}\mskip1.5mu\}}} & = & \ensuremath{\De{ \Varid{e}_{\mathrm{1}};\,\!\ldots\!\;\Varid{e}_{\Varid{n}}}} \\
  \ensuremath{\De{ \textbf{option} \{\mskip1.5mu \mskip1.5mu\}}} & = & \ensuremath{\textbf{false}} \\
  \ensuremath{\De{ \textbf{option} \{\mskip1.5mu \Varid{e}\mskip1.5mu\}}} & = & \ensuremath{\textbf{if}\;(\Varid{x}\mathrel{\coloneqq}\Varid{e})\;\textbf{then}\;\textbf{truth}\;(\Varid{x})\;\textbf{else}\;\textbf{false}} & \text{\ensuremath{\Varid{x}} fresh} \\
\end{array}
$$
{\em Note}: The desugaring assumes a scope correct Verse program.
This means that there can be no variable name clashes in \ensuremath{\mathbf{let}} and \ensuremath{\textbf{do}}.
A more robust desugaring could \textalpha-convert such clashes.
\caption{Desugaring of expressions (2)} \label{fig:desugar-expr2}
\end{figurebox}


\begin{figurebox}
\begin{hscode}\SaveRestoreHook
\column{B}{@{}>{\hspre}l<{\hspost}@{}}%
\column{8}{@{}>{\hspre}l<{\hspost}@{}}%
\column{E}{@{}>{\hspre}l<{\hspost}@{}}%
\>[B]{}\textbf{any}{}\<[8]%
\>[8]{}\mathrel{=}(\Varid{x}\Rightarrow \Varid{x}){}\<[E]%
\\
\>[B]{}\textbf{int}{}\<[8]%
\>[8]{}\mathrel{=}(\Varid{x}\Rightarrow \textbf{isInt}(\Varid{x});\,\Varid{x}){}\<[E]%
\\
\>[B]{}\textbf{nat}{}\<[8]%
\>[8]{}\mathrel{=}(\Varid{x}\Rightarrow \textbf{isInt}(\Varid{x});\,\Varid{x}\mathbin{>}(\mathbin{-}\mathrm{1});\,\Varid{x}){}\<[E]%
\\
\>[B]{}\textbf{false}{}\<[8]%
\>[8]{}\mathrel{=}\Varid{array}\;(){}\<[E]%
\\
\>[B]{}\textbf{arrow}{}\<[8]%
\>[8]{}\mathrel{=}(\Varid{ft}\Rightarrow \Varid{f}\Rightarrow \Varid{y}\Rightarrow (\Varid{ft}(\Varid{y}))((\Varid{f}(\Varid{y})))){}\<[E]%
\ColumnHook
\end{hscode}\resethooks

\caption{Core prelude} \label{fig:core-prelude}
\end{figurebox}


\begin{figurebox}

$$
  \begin{array}{rcll}
    \ensuremath{\Cr{ \cdot }}   & \multicolumn{3}{l}{ \text{ converts the limited Source Verse to Core.} } \\
    \ensuremath{\Cr{ \Varid{k}}} & = & \ensuremath{\Varid{k}} \\
    \ensuremath{\Cr{ \Varid{x}}} & = & \ensuremath{\Varid{x}} \\
    \ensuremath{\Cr{ \Varid{array}\;(\Varid{e}_{\mathrm{1}},\!\ldots\!,\Varid{e}_{\Varid{n}})}} & = &
    \begin{array}[t]{@{}l}
    \ensuremath{\textbf{def}\;\Varid{x}_{\mathrm{1}},\!\ldots\!,\Varid{x}_{\Varid{n}}\;\textbf{in}\\\{\Varid{x}_{\mathrm{1}}\mathrel{=}\Varid{e}_{\mathrm{1}};\,\!\ldots\!;\,\Varid{x}_{\Varid{n}}\mathrel{=}\Varid{e}_{\Varid{n}};\,\langle\Varid{v}_{\mathrm{1}},\!\ldots\!,\Varid{v}_{\Varid{n}}\rangle\}} \\
    \end{array}
	& \text{\ensuremath{\Varid{x}_{\Varid{i}}} fresh} \\
    \ensuremath{\Cr{ \Varid{e}_{\mathrm{1}};\,\Varid{e}_{\mathrm{2}}}} & = & \ensuremath{\Cr{ \Varid{e}_{\mathrm{1}}};\,\Cr{ \Varid{e}_{\mathrm{2}}}} \\
    \ensuremath{\Cr{ \Varid{e}_{\mathrm{1}}\mathrel{=}\Varid{e}_{\mathrm{2}}}} & = & \ensuremath{\Cr{ \Varid{e}_{\mathrm{1}}}\mathrel{=}\Cr{ \Varid{e}_{\mathrm{2}}}} \\
    \ensuremath{\Cr{ \Varid{x}\mathrel{\coloneqq}\mathbin{:}\textbf{any}}} & = & \ensuremath{\Varid{x}} \\
    \ensuremath{\Cr{ \Varid{x}\mathrel{\coloneqq}\Varid{e}}} & = & \ensuremath{\Varid{x}\mathrel{=}\Cr{ \Varid{e}}} \\
    \ensuremath{\Cr{ \Varid{e}_{\mathrm{1}} [\mskip1.5mu \Varid{e}_{\mathrm{2}}\mskip1.5mu]}} & = & \ensuremath{\exists\Varid{x},\Varid{y}.\,\Varid{x}\mathrel{=}\Varid{e}_{\mathrm{1}};\,\Varid{y}\mathrel{=}\Varid{e}_{\mathrm{2}};\,\Cr{ \Varid{e}_{\mathrm{1}}}(\Cr{ \Varid{e}_{\mathrm{2}}})} & \text{\ensuremath{\Varid{x},\Varid{y}} fresh} \\
    \ensuremath{\Cr{ \Varid{e}_{\mathrm{1}} (\Varid{e}_{\mathrm{2}})}} & = & \ensuremath{\textbf{succeeds}\,\{\Cr{ \Varid{e}_{\mathrm{1}} [\mskip1.5mu \Varid{e}_{\mathrm{2}}\mskip1.5mu]}\}} \\
    \ensuremath{\Cr{ \Varid{e}_{\mathrm{1}}\hspace{0.2em}\mathop{\rule[-0.1em]{0.15em}{0.75em}}\hspace{0.2em}\Varid{e}_{\mathrm{2}}}} & = & \ensuremath{\Bi{ \Varid{e}_{\mathrm{1}}}\hspace{0.2em}\mathop{\rule[-0.1em]{0.15em}{0.75em}}\hspace{0.2em}\Bi{ \Varid{e}_{\mathrm{2}}}} \\
    \ensuremath{\Cr{ \textbf{if}\;\Varid{e}_{\mathrm{1}}\;\textbf{then}\;\Varid{e}_{\mathrm{2}}\;\textbf{else}\;\Varid{e}_{\mathrm{3}}}} & = &
    \begin{array}[t]{@{}l}
    \ensuremath{\textbf{def}\;\Varid{x}\;\textbf{in}\\\{\Varid{x}\mathrel{=}\textbf{one}\{\Bi{ \Varid{e}_{\mathrm{1}};\,\Th{ \Varid{e}_{\mathrm{2}}}}\hspace{0.2em}\mathop{\rule[-0.1em]{0.15em}{0.75em}}\hspace{0.2em}\Bi{ \Th{ \Varid{e}_{\mathrm{3}}}}\};\,\Varid{x}(())\}}
    \end{array}
    & \text{\ensuremath{\Varid{x}} fresh} \\
    \ensuremath{\Cr{ \textbf{for}\;\Varid{e}_{\mathrm{1}}\;\textbf{in}\;\Varid{e}_{\mathrm{2}}}} & = & \ensuremath{\textbf{all}\{\Bi{ \Varid{e}_{\mathrm{1}};\,\Th{ \Varid{e}_{\mathrm{2}}}}\}} \\
    \ensuremath{\Cr{ \textbf{fun}(\Varid{x}\mathbin{:}\textbf{any})\Varid{e}}} & = & \ensuremath{\Varid{x}\Rightarrow \Bi{ \Varid{e}}} \\
%%    |Cr (type^{e})| & = & |x => Bi (x = e; x)| & \text{|x| fresh} \\
    \\
    \ensuremath{\Th{ \cdot }} & \multicolumn{3}{l}{ \text{creates a thunk, i.e., dummy function.} } \\
    \ensuremath{\Th{ \Varid{e}}} & = & \ensuremath{\textbf{fun}(\Varid{x}\mathbin{:}\textbf{any})\Varid{e}} & \text{\ensuremath{\Varid{x}} fresh} \\
    \\
    \ensuremath{\Bi{ \cdot }} & \multicolumn{3}{l}{ \text{inserts \ensuremath{\exists{}} for all local bindings.} } \\
    \ensuremath{\Bi{ \Varid{e}}} & = & \ensuremath{\exists\Vr{ \Varid{e}}.\,\Cr{ \Varid{e}}} \\
    \\
    \ensuremath{\Vr{ \cdot }} & \multicolumn{3}{l}{ \text{finds all locally defined variables.} } \\
    \ensuremath{\Vr{ \Varid{k}}} & = & \ensuremath{\epsilon} \\
    \ensuremath{\Vr{ \Varid{x}}} & = & \ensuremath{\epsilon} \\
    \ensuremath{\Vr{ \Varid{array}\;(\Varid{e}_{\mathrm{1}},\!\ldots\!,\Varid{e}_{\Varid{n}})}} & = & \ensuremath{\Vr{ \Varid{e}_{\mathrm{1}}},\!\ldots\!,\Vr{ \Varid{e}_{\Varid{n}}}} \\
    \ensuremath{\Vr{ \Varid{e}_{\mathrm{1}};\,\Varid{e}_{\mathrm{2}}}} & = & \ensuremath{\Vr{ \Varid{e}_{\mathrm{1}}},\Vr{ \Varid{e}_{\mathrm{2}}}} \\
    \ensuremath{\Vr{ \Varid{e}_{\mathrm{1}}\mathrel{=}\Varid{e}_{\mathrm{2}}}} & = & \ensuremath{\Vr{ \Varid{e}_{\mathrm{1}}},\Vr{ \Varid{e}_{\mathrm{2}}}} \\
    \ensuremath{\Vr{ \Varid{x}\mathrel{\coloneqq}\Varid{e}}} & = & \ensuremath{\Varid{x},\Vr{ \Varid{e}}} \\
    \ensuremath{\Vr{ \Varid{e}_{\mathrm{1}} [\mskip1.5mu \Varid{e}_{\mathrm{2}}\mskip1.5mu]}} & = & \ensuremath{\Vr{ \Varid{e}_{\mathrm{1}}},\Vr{ \Varid{e}_{\mathrm{2}}}} \\
    \ensuremath{\Vr{ \Varid{e}_{\mathrm{1}} (\Varid{e}_{\mathrm{2}})}} & = & \ensuremath{\Vr{ \Varid{e}_{\mathrm{1}}},\Vr{ \Varid{e}_{\mathrm{2}}}} \\
    \ensuremath{\Vr{ \Varid{e}_{\mathrm{1}}\hspace{0.2em}\mathop{\rule[-0.1em]{0.15em}{0.75em}}\hspace{0.2em}\Varid{e}_{\mathrm{2}}}} & = & \ensuremath{\epsilon} \\
    \ensuremath{\Vr{ \textbf{if}\;\Varid{e}_{\mathrm{1}}\;\textbf{then}\;\Varid{e}_{\mathrm{2}}\;\textbf{else}\;\Varid{e}_{\mathrm{3}}}} & = & \ensuremath{\epsilon} \\
    \ensuremath{\Vr{ \textbf{for}\;\Varid{e}_{\mathrm{1}}\;\textbf{in}\;\Varid{e}_{\mathrm{2}}}} & = & \ensuremath{\epsilon} \\
    \ensuremath{\Vr{ \textbf{fun}(\Varid{x}\mathbin{:}\textbf{any})\Varid{e}}} & = & \ensuremath{\epsilon} \\
    \ensuremath{\Vr{ \textbf{type} \{\mskip1.5mu \Varid{e}\mskip1.5mu\}}} & = & \ensuremath{\epsilon} \\
  \end{array}
  $$
\caption{Core conversion} \label{fig:core-conversion}
\end{figurebox}


\begin{figurebox}
$$
\begin{array}{l}
\begin{array}{rclll}
    \multicolumn{5}{l}{ \text{Reduction rules for \ensuremath{\textbf{split}} } } \\
    \rulename{split-fail} & \ensuremath{\textbf{split}(\textbf{fail})\{\Varid{f},\Varid{g}\}} & \movesto & \ensuremath{\Varid{f}(())} \\
    \rulename{split-value} & \ensuremath{\textbf{split}(\Varid{v})\{\Varid{f},\Varid{g}\}} & \movesto & \ensuremath{\exists\Varid{x}.\,\Varid{x}\mathrel{=}\Varid{g}(\Varid{v});\,\Varid{x}((\anonymous \Rightarrow \textbf{fail}))} & \text{\ensuremath{\Varid{x}} fresh} \\
    \rulename{split-choice} & \ensuremath{\textbf{split}(\Varid{v}\hspace{0.2em}\mathop{\rule[-0.1em]{0.15em}{0.75em}}\hspace{0.2em}\Varid{e})\{\Varid{f},\Varid{g}\}} & \movesto & \ensuremath{\exists\Varid{x}.\,\Varid{x}\mathrel{=}\Varid{g}(\Varid{v});\,\Varid{x}((\anonymous \Rightarrow \Varid{e}))} & \text{\ensuremath{\Varid{x}} fresh} \\
    \rulename{split-wrong} & \ensuremath{\textbf{split}(\textbf{wrong})\{\Varid{f},\Varid{g}\}} & \movesto & \ensuremath{\textbf{wrong}} \\
    \rulename{split-wrong} & \ensuremath{\textbf{split}(\textbf{wrong}\hspace{0.2em}\mathop{\rule[-0.1em]{0.15em}{0.75em}}\hspace{0.2em}\Varid{e})\{\Varid{f},\Varid{g}\}} & \movesto & \ensuremath{\textbf{wrong}} \\
\end{array} \\
\\
\begin{array}{rclll}
    \multicolumn{5}{l}{ \text{Definition of \ensuremath{\textbf{one}}, \ensuremath{\textbf{all}}, \ensuremath{\textbf{succeeds}}, and \ensuremath{\textbf{decides}} in terms of \ensuremath{\textbf{split}}. } } \\
    \ensuremath{\textbf{one}\{\Varid{e}\}} & = & \ensuremath{\textbf{split}(\Varid{e})\{\anonymous \Rightarrow \textbf{fail},\Varid{x}\Rightarrow \anonymous \Rightarrow \Varid{x}\}} & \text{\ensuremath{\Varid{x}} fresh} \\
    \ensuremath{\textbf{all}\{\Varid{e}\}} & = &
    \begin{array}[t]{@{}l}
      \ensuremath{\exists{}\;\Varid{f},\Varid{g}\;\textbf{in}\;\{\mskip1.5mu } \\
      \quad \ensuremath{\Varid{f}\mathrel{\coloneqq}\anonymous \Rightarrow ();\,} \\
      \quad \ensuremath{\Varid{g}\mathrel{\coloneqq}\Varid{x}\Rightarrow \Varid{y}\Rightarrow \textbf{cons}(\Varid{x}(()),\textbf{split}(\Varid{y}(()))\{\Varid{f},\Varid{g}\});\,} \\
      \quad \ensuremath{\textbf{split}(\Varid{e})\{\Varid{f},\Varid{g}\}\mskip1.5mu\}}
    \end{array}  & \text{\ensuremath{\Varid{x},\Varid{y}} fresh} \\
    \ensuremath{\textbf{succeeds}\,\{\Varid{e}\}} & = & \ensuremath{\textbf{split} \{\mskip1.5mu \Varid{e},\anonymous \Rightarrow \textbf{wrong},} \\
    \quad        &   & \quad \ensuremath{\Varid{x}\Rightarrow \Varid{y}\Rightarrow \textbf{split}(\Varid{y}(()))\{\anonymous \Rightarrow \Varid{x},\anonymous \Rightarrow \anonymous \Rightarrow \textbf{wrong}\}\mskip1.5mu\}} & \text{\ensuremath{\Varid{x},\Varid{y}} fresh}\\
    \ensuremath{\textbf{decides}\,\{\Varid{e}\}} & = & \ensuremath{\textbf{split} \{\mskip1.5mu \Varid{e},\anonymous \Rightarrow \textbf{fail},} \\
    \quad        &   & \quad \ensuremath{\Varid{x}\Rightarrow \Varid{y}\Rightarrow \textbf{split}(\Varid{y}(()))\{\anonymous \Rightarrow \Varid{x},\anonymous \Rightarrow \anonymous \Rightarrow \textbf{wrong}\}\mskip1.5mu\}} & \text{\ensuremath{\Varid{x},\Varid{y}} fresh}\\
\end{array}
\end{array}
$$
\lennart{Does \ensuremath{\Varid{f}} need to be a thunk?  Does second argument to \ensuremath{\Varid{g}} need to be a thunk? Yes, or it may \ensuremath{\textbf{fail}} in the wrong place.}
\caption{Alternative: Use \ensuremath{\textbf{split}} instead of \ensuremath{\textbf{one}}, \ensuremath{\textbf{all}}, \ensuremath{\textbf{succeeds}}, and \ensuremath{\textbf{decides}}.} \label{fig:split}
\end{figurebox}



\end{document}
